\begin{abstract}
Regenerating organisms restore their bodies after severe damage. Whether they
restore what they \emph{knew} is a different question, and a sharper one when the
knowing is distributed: computation that lives nowhere in particular cannot be
deleted region by region, only perturbed. We study this in neural cellular
automata, where one shared local update rule grows a body from a single cell and
regrows it after tissue loss, so damage removes computation rather than a
component. Two such organisms play a Lewis referential game --- a sender emits a
discrete symbol through a parameter-free consensus readout over its own living
cells, and a receiver that observes nothing else must identify the referent. We
then damage the sender while holding the receiver frozen, so an unchanged partner
acts as a fixed semantic reference frame and any lost accuracy is the sender's
protocol drifting from the meaning its partner reads. Our thesis is that
regeneration restores the body and the channel capacity but not the convention.
%
Because no definition of semantic recovery exists to borrow, we contribute one:
four operationally distinct recoveries, a decomposition separating whether the
system still speaks \emph{a} language from whether it speaks the \emph{same} one,
and an exact zero-cost control separating genuine drift from the mechanical
effect of averaging over fewer cells. We prove that under full per-cell
consensus, cell deletion cannot flip the emitted symbol, converting the central
confound from an argument into a theorem. Empirically the substrate grows and
regenerates reliably (IoU 0.9898 at the training horizon, 0.9993 at four times
it; 256 of 256 damage trials recover, median 13 steps), and we measure a limit
that reshaped the design: a morphology-trained automaton transports an injected
signal only about two cells before its propagation front stalls, against a body
of radius nine. Communication across such a body must be trained for, not
assumed. \gated{protocol emergence and the dissociation experiment}
\end{abstract}
